\documentclass[10pt,a4paper]{amsart}
\usepackage[margin=19mm]{geometry}
\usepackage{amsmath,amssymb,mathtools}
\usepackage[scr=rsfs,cal=euler]{mathalfa}
\usepackage{xfrac}
\usepackage[unicode,hidelinks,bookmarks=false]{hyperref}
\newcommand{\ie}{i.e.\ }
\newcommand{\cf}{cf.\ }
\newcommand{\resp}{respectively\ }
\newcommand{\Subsection}{\subsection}
\providecommand{\nobreakdash}{\nobreak\hbox{-}}
\setlength{\emergencystretch}{2em}
\multlinegap0pt
\usepackage{bm}
\usepackage{bbm}
\usepackage{dsfont}
\usepackage{xspace}
\usepackage{cancel}
\usepackage{mathtools}
\usepackage{amsthm}
\makeatletter
\@ifundefined{theorem}{\newtheorem{theorem}{Theorem}}{}
\@ifundefined{lemma}{\newtheorem{lemma}[theorem]{Lemma}}{}
\makeatother
\title[On the Structure of 3D Queen Domination]{On the Structure of 3D Queen Domination}
\author{Mahesh Ramani}
\date{Source: arXiv:2604.03793v1 (CC BY 4.0)}
\begin{document}
\maketitle

\noindent\emph{Source and licence.} On the Structure of 3D Queen Domination. Version arXiv:2604.03793v1 (CC BY 4.0), distributed under the Creative Commons Attribution 4.0 International licence, as recorded in the arXiv licence metadata for this exact version and shown on the abstract page https://arxiv.org/abs/2604.03793v1. Full rights notice: licenses/arxiv-2604.03793-CC-BY-4.0.txt.\par\smallskip

\noindent\textit{Editorial scope.} Counted unit: the Theorem whose source label is \texttt{thm:orbit} in the article (source0.tex lines 125--144, source label \texttt{thm:orbit}), delivered together with the article's own complete original proof of it (source0.tex lines 146--202, one \texttt{proof} environment of 57 lines). No mathematics is written, repaired or re-derived: statement and proof are the article's own text line for line. Required context retained uncounted: lem:parity (lines 97--109). Every definition, display and label of those lines is retained unchanged. Omitted from this package and not redistributed: 1--96, 110--124, 145--145, 203--425. Nothing inside a retained excerpt is omitted or altered: every retained body is delivered exactly as the article's lines are written, so no omission receipt of any kind accompanies this package. Citations: the retained excerpts carry no citation command at all, so no bibliography entry of the article is transcribed and no reference is redistributed. Reference adaptation: none was needed and none was made; every reference inside the delivered unit resolves inside it, so no unresolved reference or citation is produced. Printed slips kept literal: no printed word, symbol, hypothesis or number of the counted statement or of its proof was altered. Layout and class: the article's own class is \texttt{article} with packages including amsmath, amsthm, amssymb, hyperref, booktabs, geometry, makecell; the wrapper is the lane's pinned \texttt{amsart} editorial wrapper at 10pt on A4, which restates the theorem-like environments the retained text needs and adds the article's own macro definitions that the retained text needs (none), with extra wrapper packages (bm, bbm, dsfont, xspace, cancel, mathtools). The delivered numbering is the unit's own. Assets: no retained span contains a figure environment, a graphics-inclusion command, a PDF-page inclusion, a table or a TikZ picture; no image file is copied into this package, the package declares no asset dependency and no image file is redistributed with it. Licence: version arXiv:2604.03793v1 of this article is distributed under the Creative Commons Attribution 4.0 International licence, as recorded in the arXiv licence metadata for this exact version and shown on the abstract page https://arxiv.org/abs/2604.03793v1. A full rights notice is delivered at licenses/arxiv-2604.03793-CC-BY-4.0.txt. Characters: every retained excerpt is pure ASCII, so the delivered engine, XeLaTeX with its default Latin Modern text font, reports no missing character.
\begin{lemma}
\label{lem:parity}
Let $M \geq 1$ be an integer and let $f(a,c) = |a-c| + |a+c-M|$ for integers
$0 \leq a, c \leq M$. Then $f(a,c) \equiv M \pmod{2}$ for all $a, c$. Hence
$f(a,c) \geq M \bmod 2$. If $M$ is even, equality holds if and only if
$a = c = M/2$. If $M$ is odd, equality holds exactly at
\[
(a,c) \in \left\{\left(\frac{M-1}{2},\frac{M-1}{2}\right),
\left(\frac{M-1}{2},\frac{M+1}{2}\right),
\left(\frac{M+1}{2},\frac{M-1}{2}\right),
\left(\frac{M+1}{2},\frac{M+1}{2}\right)\right\}.
\]
\end{lemma}

\begin{theorem}
\label{thm:orbit}
Let $n \geq 4$ and $m = n-2$.
\begin{enumerate}
  \item[\textnormal{(i)}] If $q$ is a corner vertex, then $\kappa(q) = m$.
  \item[\textnormal{(ii)}] If $q$ lies on an edge but not at a corner, then
        $\kappa(q) = 2m-1$.
  \item[\textnormal{(iii)}] If $q$ lies on a face but not on any edge, then
        \[
          \kappa(q) \;\leq\; 5m - 4 - \varepsilon_m,
          \qquad \varepsilon_m = (m-1)\bmod 2,
        \]
        with equality when the face centre is a lattice point, which occurs
        precisely when $m$ is odd.
  \item[\textnormal{(iv)}] If $q \in C$, then $\kappa(q) \leq 13m-12$, with
        equality when $q$ is at the centre of the board.
\end{enumerate}
In particular, $\max_{q\in\partial[n]^3}\kappa(q) < \max_{q\in C}\kappa(q)$
for all $n \geq 4$.
\end{theorem}

\begin{proof}
Since the 13 queen-line directions through any cell are mutually non-parallel,
two distinct queen lines through the same cell intersect only at that cell.

\smallskip\noindent\textbf{(iv).}
For $q \in C$, each of the 13 queen directions determines a segment; its
intersection with $C$ contains at most $m$ cells (including $q$). Since the
segments share only $q$,
\[
  \kappa(q) \;\leq\; 13(m-1)+1 \;=\; 13m-12.
\]
A cell at the centre of $C$ attains all 13 full segments, so equality holds.

\smallskip\noindent\textbf{(i).}
At $q = (0,0,0)$, a direction $(d_x,d_y,d_z)$ reaches $C = \{1,\ldots,m\}^3$
only if $d_x = d_y = d_z = 1$. That space diagonal visits exactly $m$ core
cells, so $\kappa(q) = m$.

\smallskip\noindent\textbf{(ii).}
At $q = (0,0,z)$ with $1 \leq z \leq m$, the directions reaching $C$ must have
$d_x = d_y = 1$. These are $(1,1,0)$, $(1,1,1)$, $(1,1,-1)$, with
core-segment lengths $m$, $m-z$, $z-1$ respectively. Thus
\[
  \kappa(q) \;=\; m + (m-z) + (z-1) \;=\; 2m-1.
\]

\smallskip\noindent\textbf{(iii).}
At $q = (0,y,z)$ with  $1 \leq y,z \leq m$, set $a = y-1$, $b = m-y$,
$c = z-1$, $d = m-z$, $M = m-1$, so $a+b = c+d = M$. All nine directions with
$d_x = 1$ reach $C$; their core-segment lengths are
\[
  m,\quad a,\quad b,\quad c,\quad d,\quad
  \min(a,c),\quad \min(a,d),\quad \min(b,c),\quad \min(b,d).
\]
Using $b = M-a$, $d = M-c$, and $\min(u,v) = \tfrac{1}{2}(u+v-|u-v|)$,
\[
  \min(a,c)+\min(a,d)+\min(b,c)+\min(b,d) \;=\; 2M - |a-c| - |a+c-M|.
\]
Therefore
\[
  \kappa(q) \;=\; 5m-4 - f(a,c), \quad f(a,c) = |a-c|+|a+c-M|.
\]
By Lemma~\ref{lem:parity}, $f(a,c) \geq M \bmod 2 = (m-1) \bmod 2 =
\varepsilon_m$, giving $\kappa(q) \leq 5m-4-\varepsilon_m$. When $M$ is even
(i.e., $m$ is odd), Lemma~\ref{lem:parity} gives a unique minimiser $a = c =
M/2$, corresponding to $y = z = (m+1)/2$, the unique lattice centre of the
face. When $M$ is odd ($m$ even), the face has no lattice centre and the
minimum $\varepsilon_m = 1$ is attained at the four lattice points closest to
the geometric centre.

\smallskip\noindent\textbf{Comparison.}
For $m \geq 2$, the boundary coverage is at most $5m - 4 - \varepsilon_m$. By (iv), the interior maximum is $13m - 12$ when $m$ is odd. When $m$ is even, evaluating any of the eight central-most cells in $C$ yields a coverage of exactly $13m - 18$. In either case, $\max_{q\in C}\kappa(q) \geq 13m - 18$. We thus have the separation:
\[
  m \;<\; 2m-1 \;\leq\; 5m-5 \;\leq\; 5m-4-\varepsilon_m \;<\; 13m-18 \;\leq\; \max_{q\in C}\kappa(q),
\]
so $\max_{q\in\partial[n]^3}\kappa(q) < \max_{q\in C}\kappa(q)$.
\end{proof}

\end{document}
